\section*{Abstract}
		Wir zeigen, dass die Koide-Formel für Leptonmassen eng mit der geometrischen Konstante $\xi = \frac{4}{3} \times 10^{-4}$ aus der T0-Theorie verknüpft ist. Die Quantenverhältnisse $(r,p)$ der T0-Yukawa-Formel $m = r \cdot \xi^p \cdot v$ erzeugen ohne zusätzliche Parameter oder fraktale Korrekturen näherungsweise die Koide-Symmetrie: $Q_{\text{T0}} = 0.667742$ liegt $0.161\,\%$ über $\frac{2}{3}$, während die empirische Relation mit den PDG-Massen auf $|\Delta Q| \approx 0.0003\,\%$ erfüllt ist.

	
	\section{Die Koide-Formel}
	
	Die 1981 von Yoshio Koide entdeckte Relation verbindet die Massen der geladenen Leptonen:
	
	\begin{equation}
		Q = \frac{m_e + m_\mu + m_\tau}{\left( \sqrt{m_e} + \sqrt{m_\mu} + \sqrt{m_\tau} \right)^2} = \frac{2}{3}
		\label{eq:koide}
	\end{equation}
	
	Diese Formel erreicht mit den PDG-2024-Massen ($m_\tau = 1776.93 \pm 0.09$~MeV) eine experimentelle Genauigkeit von $|\Delta Q| \approx 0.0003\%$; die Unsicherheit von $m_\tau$ allein ergibt $\pm 0.0008\%$.
	
	\section{T0-Yukawa-Formel}
	
	In der T0-Theorie entstehen Teilchenmassen durch:
	
	\begin{equation}
		m = r \cdot \xi^p \cdot v
		\label{eq:t0yukawa}
	\end{equation}
	
	mit Higgs-VEV $v = 246$ GeV und $\xi = \frac{4}{3} \times 10^{-4}$.
	
	\subsection{Leptonparameter}
	
	\begin{table}[h]
		\centering
		\adjustbox{max width=\textwidth}{%
\begin{tabular}{lccc}
			\toprule
			\textbf{Lepton} & \textbf{$r$} & \textbf{$p$} & \textbf{$m$ [GeV]} \\
			\midrule
			Elektron & $\frac{4}{3}$ & $\frac{3}{2}$ & 0.000511 \\
			Myon & $\frac{16}{5}$ & $1$ & 0.1057 \\
			Tau & $\frac{25}{9}$ & $\frac{2}{3}$ & 1.7769 \\
			\bottomrule
		\end{tabular}%
}
		\caption{T0-Quantenverhältnisse der geladenen Leptonen}
	\end{table}
	
	\section{Haupttheorem}
	
	\begin{theorem}
		Die Koide-Relation $Q \approx \frac{2}{3}$ folgt näherungsweise aus den T0-Exponenten $(p_e, p_\mu, p_\tau) = \left(\frac{3}{2}, 1, \frac{2}{3}\right)$ und den zugehörigen Verhältnissen $(r_e, r_\mu, r_\tau) = \left(\frac{4}{3}, \frac{16}{5}, \frac{25}{9}\right)$; mit diesen Werten ist $Q_{\text{T0}} = 0.667742$, also $0.161\,\%$ über $\frac{2}{3}$.
	\end{theorem}
	
	\section{Beweis durch Massenverhältnisse}
	
	\subsection{Elektron zu Myon}
	
	\begin{beweis}
		\allowdisplaybreaks
		\begin{align}
			\frac{m_e}{m_\mu} &= \frac{r_e \cdot \xi^{p_e}}{r_\mu \cdot \xi^{p_\mu}} = \frac{\frac{4}{3} \cdot \xi^{3/2}}{\frac{16}{5} \cdot \xi^1} \\
			&= \frac{4}{3} \cdot \frac{5}{16} \cdot \xi^{1/2} = \frac{5}{12} \cdot \xi^{1/2} \\
			&= \frac{5}{12} \cdot \sqrt{1.333 \times 10^{-4}} \\
			&= \frac{5}{12} \cdot 0.01155 = 0.004813 \\
			&\approx \frac{1}{206.768} \quad \checkmark
		\end{align}
		
		\textbf{Experimentell:} $\frac{m_e}{m_\mu} = 0.004836$ (PDG 2024)\\
		\textbf{Abweichung:} $< 0.5\%$
	\end{beweis}
	
	\subsection{Myon zu Tau}
	
	\begin{beweis}
		\begin{align}
			\frac{m_\mu}{m_\tau} &= \frac{r_\mu \cdot \xi^{p_\mu}}{r_\tau \cdot \xi^{p_\tau}} = \frac{\frac{16}{5} \cdot \xi^1}{\frac{25}{9} \cdot \xi^{2/3}} \\
			&= \frac{16}{5} \cdot \frac{9}{25} \cdot \xi^{1/3} = \frac{144}{125} \cdot \xi^{1/3} \\
			&= 1.152 \cdot (1.333 \times 10^{-4})^{1/3} \\
			&= 1.152 \cdot 0.05105 = 0.05881 \\
			&\approx \frac{1}{17.0} \quad \checkmark
		\end{align}
		
		\textbf{Experimentell:} $\frac{m_\mu}{m_\tau} = 0.05947$ (PDG 2024)\\
		\textbf{Abweichung:} $< 1.2\%$
	\end{beweis}
	
	\subsection{Elektron zu Tau}
	
	\begin{beweis}
		\begin{align}
			\frac{m_e}{m_\tau} &= \frac{r_e \cdot \xi^{p_e}}{r_\tau \cdot \xi^{p_\tau}} = \frac{\frac{4}{3} \cdot \xi^{3/2}}{\frac{25}{9} \cdot \xi^{2/3}} \\
			&= \frac{4}{3} \cdot \frac{9}{25} \cdot \xi^{5/6} = \frac{12}{25} \cdot \xi^{5/6} \\
			&= 0.48 \cdot (1.333 \times 10^{-4})^{5/6} \\
			&= 0.48 \cdot 0.000590 = 0.000283 \\
			&\approx \frac{1}{3534} \quad \checkmark
		\end{align}
		
		\textbf{Experimentell:} $\frac{m_e}{m_\tau} = 0.0002876$ (PDG 2024)\\
		\textbf{Abweichung:} $< 1.6\%$
	\end{beweis}
	
	\section{Direkte Herleitung der Koide-Relation}
	
	\subsection{Geometrische Struktur der Exponenten}
	
	Die T0-Exponenten zeigen eine fundamentale Symmetrie:
	
	\begin{equation}
		p_e - p_\mu = \frac{3}{2} - 1 = \frac{1}{2}
	\end{equation}
	\begin{equation}
		p_\mu - p_\tau = 1 - \frac{2}{3} = \frac{1}{3}
	\end{equation}
	
	Diese erzeugen die charakteristischen $\sqrt{m}$-Abhängigkeiten der Koide-Formel.
	
	\subsection{Berechnung von $Q$}
	
	Setzen wir die T0-Massen in Gleichung \eqref{eq:koide} ein:
	
	\begin{align}
		Q &= \frac{r_e \xi^{p_e} v + r_\mu \xi^{p_\mu} v + r_\tau \xi^{p_\tau} v}{\left(\sqrt{r_e \xi^{p_e} v} + \sqrt{r_\mu \xi^{p_\mu} v} + \sqrt{r_\tau \xi^{p_\tau} v}\right)^2} \\
		&= \frac{r_e \xi^{3/2} + r_\mu \xi + r_\tau \xi^{2/3}}{\left(\sqrt{r_e} \xi^{3/4} + \sqrt{r_\mu} \xi^{1/2} + \sqrt{r_\tau} \xi^{1/3}\right)^2}
	\end{align}
	$v$ kürzt sich in Zähler und Nenner vollständig, $Q$ ist dimensionslos.
	
	Mit den numerischen Werten ($m = (0.505, 104.96, 1783.4)$~MeV):
	\begin{align}
		Q_{\text{T0}} &= 0.667742 \\
		Q_{\text{Koide}} &= \frac{2}{3} = 0.666667 \\
		\Delta Q &= +0.161\%
	\end{align}
	Mit diesen $(r,p)$ ist $Q = 2/3$ numerisch nicht exakt erfüllt; die Folgerungen gelten näherungsweise ($Q \approx 2/3$).
	
	\section{Schlüsselerkenntnis}
	
	\begin{folgerung}
		\textbf{Die Koide-Formel ist keine unabhängige Symmetrie, sondern eine direkte Manifestation von $\xi$.}
		
		\begin{itemize}
			\item Die Exponenten $(3/2, 1, 2/3)$ erzeugen die $\sqrt{m}$-Struktur
			\item Die Verhältnisse $(4/3, 16/5, 25/9)$ sind konsistent mit $Q \approx 2/3$
			\item Keine fraktalen Korrekturen nötig
			\item Keine zusätzlichen freien Parameter
			\item Die geometrische Konstante $\xi$ war implizit bereits in der Koide-Formel enthalten
		\end{itemize}
	\end{folgerung}
	
	\section{Vergleich: Empirische vs. T0-Herleitung}
	
	\begin{table}[h]
		\centering
		\adjustbox{max width=\textwidth}{%
\begin{tabular}{lcc}
			\toprule
			\textbf{Aspekt} & \textbf{Koide (1981)} & \textbf{T0-Theorie} \\
			\midrule
			Freie Parameter & 0 (empirisch) & 1 ($\xi$) \\
			Basis & Beobachtung & Geometrie \\
			Genauigkeit & $\approx 0.0003\%$ & $0.16\%$ \\
			Erklärung & Keine & $\xi$-Geometrie \\
			Vorhersagekraft & Nur Leptonen & Alle Teilchen \\
			\bottomrule
		\end{tabular}%
}
		\caption{Vergleich der Ansätze}
	\end{table}
	
	\section{Mathematische Bedeutung}
	
	Die T0-Formel zeigt, dass:
	
	\begin{equation}
		Q \approx \frac{2}{3} \iff \text{Exponenten bilden geometrische Reihe mit Basis } \xi
	\end{equation}
	
	Dies erklärt:
	\begin{enumerate}
		\item Warum $Q$ nahe bei $2/3$ liegt und nicht bei einem anderen Wert
		\item Warum die Relation für genau 3 Generationen gilt
		\item Warum Wurzeln der Massen (nicht Massen selbst) addiert werden
		\item Die Verbindung zur Higgs-Yukawa-Kopplung
	\end{enumerate}
	
	\section{Feinstrukturkonstante und Massenverhältnisse}
	
	\subsection{Direkte T0-Ableitung}
	
	Die Feinstrukturkonstante in der T0-Theorie:
	
	\begin{equation}
		\alpha = \xi \cdot \left(\frac{E_0}{1\,\text{MeV}}\right)^2 = \frac{4}{3} \times 10^{-4} \times (7.398)^2 = 0.007297
	\end{equation}
	
	mit der Energieskala $E_0$; der im folgenden Unterabschnitt gezeigte Versuch, $\alpha$ allein aus den Lepton-Massenverhältnissen zu rekonstruieren, trifft $\alpha$ nicht.
	
	\textbf{Experimentell:} $\alpha = \frac{1}{137.036} = 0.0072973525693$\\
	\textbf{Fehler:} $0.0005\%$ ($1/(\xi\cdot7.398^2) = 137.035$ gegen $137.036$)
	
	\subsection{Rekonstruktion aus Leptonmassen}
	
	\begin{beweis}
		Ein Rekonstruktionsversuch der Feinstrukturkonstante aus den Massenverhältnissen:
		
		\begin{equation}
			\alpha \propto \left(\frac{m_e}{m_\mu}\right)^{2/3} \times \left(\frac{m_\mu}{m_\tau}\right)^{1/2} \times \xi^{\text{konst}}
		\end{equation}
		
		Mit den T0-Verhältnissen:
		\begin{align}
			\alpha_{\text{rekon}} &= \left(\frac{1}{206.768}\right)^{2/3} \times \left(\frac{1}{16.818}\right)^{1/2} \times 1.089 \\
			&= 0.02860 \times 0.2438 \times 1.089 \\
			&\approx 0.00759
		\end{align}
		Das ergibt $1/\alpha = 131.7$ ($+4\,\%$); der Faktor $1.089$ ist nicht hergeleitet. Die Rekonstruktion trifft $\alpha$ damit nicht.
	\end{beweis}
	
	\textbf{Zuordnung der Exponenten:} Die T0-Exponenten-Differenzen bestimmen die $\xi$-Abhängigkeit der beiden Verhältnisse:
	\begin{itemize}
		\item $p_e - p_\mu = \frac{3}{2} - 1 = \frac{1}{2}$ gehört zu $m_e/m_\mu \propto \xi^{1/2}$
		\item $p_\mu - p_\tau = 1 - \frac{2}{3} = \frac{1}{3}$ gehört zu $m_\mu/m_\tau \propto \xi^{1/3}$
	\end{itemize}
	Die Exponenten $(2/3, 1/2)$ der Rekonstruktion stehen dazu über Kreuz: Der Exponent $1/2$ sitzt auf $m_\mu/m_\tau$, nicht auf $m_e/m_\mu$. Eine direkte Verknüpfung mit den Exponenten-Differenzen besteht daher nicht.

	
	\section{Hierarchie der $\xi$-Manifestationen}
	
	Die drei fundamentalen Konstanten entstehen aus $\xi$ auf verschiedenen ``Reinheits-Ebenen'':
	
	\subsection{Ebene 1: Massenverhältnisse (Koide-Formel)}
	
	\begin{equation}
		Q = \frac{\sum m_i}{\left(\sum \sqrt{m_i}\right)^2} \quad \text{mit} \quad m_i = r_i \xi^{p_i} v
	\end{equation}
	
	\begin{tcolorbox}[colback=green!5!white,colframe=green!75!black,title=Reinste $\xi$-Form]
		\textbf{Genauigkeit:} $\Delta Q = +0.161\%$ mit $m_i = r_i \xi^{p_i} v$; empirisch (PDG-Massen) $|\Delta Q| \approx 0.0003\%$
		
		\textbf{Struktur:}
		\begin{itemize}
			\item Nur Verhältnisse, keine Absolutskalen
			\item $\xi$ erscheint nur in Exponenten-Differenzen: $\xi^{p_i - p_j}$
			\item Higgs-VEV $v$ kürzt sich vollständig
			\item KEINE fraktalen Korrekturen nötig
		\end{itemize}
	\end{tcolorbox}
	
	\subsection{Ebene 2: Feinstrukturkonstante}
	
	\begin{equation}
		\alpha = \xi \cdot E_0^2
	\end{equation}
	
	\begin{tcolorbox}[colback=blue!5!white,colframe=blue!75!black,title=Semi-reine $\xi$-Form]
		\textbf{Genauigkeit:} $\Delta \alpha \approx 0.0005\%$
		
		\textbf{Warum sehr gut:}
		\begin{itemize}
			\item Benötigt eine Energieskala $E_0 = 7.398$ MeV, die aus den Massenverhältnissen emergent abgeleitet wird
			\item Direkte $\xi$-Kopplung
			\item Kleine Unsicherheit durch $E_0$-Kalibrierung
		\end{itemize}
	\end{tcolorbox}
	
	\subsection{Ebene 3: Gravitationskonstante}
	
	\begin{equation}
		G = \frac{\xi^2}{4m} = \frac{\xi^2}{4 \cdot \xi/2} = \frac{\xi}{2} = 6.667 \times 10^{-5} \quad \text{(in nat. Einheiten)}
	\end{equation}
	
	Mit SI-Umrechnung: $G_{\text{SI}} = G_{\text{nat}} \times 1.00115 \times 10^{-6}\,\text{m}^3\text{kg}^{-1}\text{s}^{-2}$
	
	Der zu $G_{\text{nat}} = \xi/2$ passende Faktor $K_{\text{SI}} = 1.00115\times10^{-6}$ entspricht Dok.~028. Die Korpusformel $G = \xi^2/(4m_e)\cdot C_{\text{conv}}\cdot K_{\text{frak}}$ liefert $6.6745\times10^{-11}$, also $+0.003\%$ zu CODATA. Dabei ist $C_{\text{conv}}$ die SI-Umrechnung der Ein-Anker-Kette (Dok.~383); ihr Zahlenwert $7.783\times10^{-3}$ enthält zusätzlich den kleinen Rest dieser Kette (Kettenwert $7.58\times10^{-3}$). Die Form von $G$ und seine Größenordnung folgen aus $\xi$ und dem einen Anker (numerisch geprüft in Dok.~383); mit dem Anker $m_e$ liegt $G$ bei $-2.6\,\%$, mit $v$ bei $-0.46\,\%$. Die Übereinstimmung auf $0.003\,\%$ ergibt sich erst mit diesem Rest und ist kein eigener Präzisionsbeleg.
	
	\begin{tcolorbox}[colback=yellow!5!white,colframe=orange!75!black,title=Komplexe $\xi$-Form]
		\textbf{Genauigkeit:} $\Delta G \approx +0.003\%$ (erst mit dem Rest der Ein-Anker-Kette in $C_{\text{conv}}$; kein eigener Präzisionsbeleg)
		
		\textbf{Warum schwieriger:}
		\begin{itemize}
			\item Benötigt Planck-Länge $\ell_P = 1.616 \times 10^{-35}$ m, die in direkter Beziehung zu $\xi$ steht ($\ell_P \propto \sqrt{G} \propto \sqrt{\xi}$ in natürlichen Einheiten)
			\item Komplexe SI-Einheiten-Umrechnung
			\item $G_{\exp}$ selbst hat $\sim 0.02\%$ Messunsicherheit
			\item Dimensionale Faktoren: $[E^{-1}] \to [E^{-2}] \to [\text{m}^3\text{kg}^{-1}\text{s}^{-2}]$
		\end{itemize}
	\end{tcolorbox}
	
	\section{Warum keine fraktalen Korrekturen?}
	
	\subsection{Verhältnis-Geometrie vs. Absolute Skalen}
	
	\begin{theorem}
		\textbf{Verhältnis-Invarianz der Koide-Formel}
		
		Die Koide-Formel arbeitet ausschließ lich mit Massenverhältnissen:
		\begin{equation}
			Q = \frac{m_e + m_\mu + m_\tau}{(\sqrt{m_e} + \sqrt{m_\mu} + \sqrt{m_\tau})^2}
		\end{equation}
		
		Da alle Massen $m_i = r_i \xi^{p_i} v$ sind, kürzen sich die $\xi$-Faktoren teilweise:
		\begin{equation}
			Q \propto \frac{\xi^{p_1} + \xi^{p_2} + \xi^{p_3}}{(\xi^{p_1/2} + \xi^{p_2/2} + \xi^{p_3/2})^2}
		\end{equation}
		
		Das Ergebnis hängt nur von den Exponenten-Differenzen ab:
		\begin{equation}
			\Delta p_{12} = p_1 - p_2, \quad \Delta p_{23} = p_2 - p_3
		\end{equation}
	\end{theorem}
	
	\subsection{Fraktale Korrekturen nur bei absoluten Skalen}
	
	\begin{table}[h]
		\centering
		\adjustbox{max width=\textwidth}{%
\begin{tabular}{lcc}
			\toprule
			\textbf{Konstante} & \textbf{Typ} & \textbf{Fraktale Korrektur?} \\
			\midrule
			$Q$ (Koide) & Verhältnis & \textbf{NEIN} \\
			$m_p/m_e$ & Verhältnis & \textbf{NEIN} \\
			$\alpha$ & Absolut mit Skala & \textbf{MINIMAL} \\
			$G$ & Absolut mit SI & \textbf{JA} \\
			\bottomrule
		\end{tabular}%
}
		\caption{Notwendigkeit fraktaler Korrekturen}
	\end{table}
	
	\section{Vereinigte Theorie der Fundamentalkonstanten}
	
	\begin{folgerung}
		\textbf{Alle drei fundamentalen Konstanten entstehen aus $\xi$:}
		
		\begin{align}
			\text{Koide: } & Q = f_1(\xi^{p_i - p_j}) \approx \frac{2}{3} \quad &&\text{(Fehler: } 0.16\%) \\
			\text{Feinstruktur: } & \alpha = \xi \cdot E_0^2 = \frac{1}{137.036} \quad &&\text{(Fehler: } 0.0005\%) \\
			\text{Gravitation: } & G = f_2(\xi, \ell_P) = 6.674 \times 10^{-11} \quad &&\text{(Fehler: } 0.003\%)
		\end{align}
		
		Die $0.003\%$ bei $G$ ergeben sich erst mit dem Rest der Ein-Anker-Kette in $C_{\text{conv}}$ (Ebene~3) und sind kein eigener Präzisionsbeleg.
		
		Die unterschiedlichen Genauigkeiten reflektieren die Komplexität der $\xi$-Manifestation.
	\end{folgerung}
	
	\subsection{Fundamentale Beziehung}
	
	Die T0-Theorie zeigt eine tiefe Verbindung:
	
	\begin{equation}
		\boxed{\xi \xrightarrow{\text{Verhältnisse}} Q \approx \frac{2}{3} \xrightarrow{\text{Skala}} \alpha \xrightarrow{\text{SI-Einheiten}} G}
	\end{equation}
	
	Jede Ebene fügt eine Komplexitätsschicht hinzu:
	\begin{itemize}
		\item \textbf{Koide:} Reine Geometrie
		\item \textbf{$\alpha$:} Geometrie + Energieskala
		\item \textbf{$G$:} Geometrie + Energieskala + Raum-Zeit-Metrik
	\end{itemize}
	
	\section{Fazit}
	
	\begin{theorem}
		\textbf{Die Koide-Formel ist die reinste $\xi$-Manifestation.}
		
		Die 1981 empirisch entdeckte Symmetrie enthielt bereits die fundamentale geometrische Konstante $\xi = \frac{4}{3} \times 10^{-4}$, ohne dass dies erkannt wurde. Die T0-Theorie zeigt:
		
		\begin{enumerate}
			\item Koide-Formel ist näherungsweise eine versteckte $\xi$-Relation ($Q_{\text{T0}}$ liegt $0.161\,\%$ über $2/3$)
			\item Feinstrukturkonstante folgt als $\alpha = \xi \cdot E_0^2$; die Rekonstruktion allein aus den Lepton-Massenverhältnissen trifft $\alpha$ nicht
			\item Gravitationskonstante ist die direkteste $\xi$-Manifestation: $G \propto \xi$
			\item Massenverhältnisse benötigen KEINE fraktalen Korrekturen
			\item Die Hierarchie $Q \to \alpha \to G$ zeigt zunehmende Komplexität
			\item Erweiterungen zu Neutrinos und Hadronen verstärken die Universalität
		\end{enumerate}
	\end{theorem}
	
	\vspace{1cm}
	
	\noindent\textbf{Historische Ironie:} Koide entdeckte 1981 eine Relation, die $\xi$ bereits enthielt, aber erst 40 Jahre später wird die geometrische Grundlage sichtbar. Die hohe Genauigkeit der Koide-Formel ($|\Delta Q| \approx 0.0003\%$ mit PDG-Massen) ist kein Zufall, sondern die Konsequenz ihrer verhältnisbasierten Natur.
	
	\begin{thebibliography}{99}
		
		\bibitem{Koide1982}
		Y. Koide, ``A fermion-boson composite model of quarks and leptons'', \textit{Phys. Lett. B} \textbf{120} (1983) 161.
		
		\bibitem{PDG2024}
		Particle Data Group, ``Review of Particle Physics'', \textit{Phys. Rev. D} \textbf{110} (2024) 030001. 
		\url{https://pdg.lbl.gov/2024/}
		
		\bibitem{T0Grundlagen}
		J. Pascher, ``T0-Theorie: Grundlagen des Zeit-Masse-Dualitäts-Frameworks'' (2024).
		
		\bibitem{T0Feinstruktur}
		J. Pascher, ``T0-Theorie: Ableitung der Feinstrukturkonstante aus $\xi$'' (2024).
		
		\bibitem{T0Gravitation}
		J. Pascher, ``T0-Theorie: Geometrische Herleitung der Gravitationskonstante'' (2024).
		
		\bibitem{T0Teilchenmassen}
		J. Pascher, ``T0-Theorie: Systematische Berechnung der Teilchenmassen'' (2024).
		
		\bibitem{T0SI}
		J. Pascher, ``T0-Theorie: SI-Reform 2019 als $\xi$-Kalibrierung'' (2024).
		
		\bibitem{T0Verhaeltnis}
		J. Pascher, ``T0-Theorie: Verhältnisse vs. absolute Werte -- Fraktale Korrekturen'' (2024).
		
		\bibitem{T0MuonG2}
		J. Pascher, ``T0-Theorie: Anomale magnetische Momente und Muon g-2'' (2024).
		
		\bibitem{T0QFT}
		J. Pascher, ``T0-Theorie: Quantenfeldtheorie und Relativitätstheorie'' (2024).
		
		\bibitem{T0Bibliographie}
		J. Pascher, ``T0-Theorie: Vollständige Bibliographie (131+ Dokumente)'' (2024).
		
		\bibitem{T0GitHub}
		J. Pascher, ``T0-Time-Mass-Duality: Complete Repository'', GitHub (2024). 
		\url{https://doi.org/10.5281/zenodo.20474821}
		
		\bibitem{T0Release}
		J. Pascher, ``T0-QFT-ML v2.0: Machine Learning Derived Extensions'', GitHub Release v1.8 (2025).
		
		\bibitem{Feynman1985}
		R. P. Feynman, ``QED: The Strange Theory of Light and Matter'', Princeton University Press (1985).
		
		\bibitem{Sommerfeld1916}
		A. Sommerfeld, ``Zur Quantentheorie der Spektrallinien'', \textit{Ann. d. Phys.} \textbf{51} (1916) 1-94.
		
		\bibitem{Dirac1937}
		P. A. M. Dirac, ``The cosmological constants'', \textit{Nature} \textbf{139} (1937) 323.
		
		\bibitem{Brannen2005}
		C. P. Brannen, ``The Lepton Masses'', \textit{arXiv:hep-ph/0501382} (2005).
		\url{https://brannenworks.com/MASSES2.pdf}
		
		\bibitem{Brannen2007}
		C. P. Brannen, ``Koide mass equations for hadrons'', \textit{arXiv:0704.1206} (2007).
		\url{http://www.brannenworks.com/koidehadrons.pdf}
		
		\bibitem{PhaseVectors2025}
		Anonymous, ``The Koide Relation and Lepton Mass Hierarchy from Phase Vectors'', \textit{rxiv.org} (2025).
		\url{https://rxiv.org/pdf/2507.0040v1.pdf}
		
		\bibitem{KoideReview2005}
		A. Rivero and A. Gsponer, ``The strange formula of Dr. Koide'', \textit{arXiv:hep-ph/0505220} (2005).
		\url{https://arxiv.org/pdf/hep-ph/0505220}
		
	\end{thebibliography}
	
\section*{Genauigkeitsangaben zur Koide-Relation}

Für die Koide-Relation werden drei Zahlen verwendet, die verschiedene Dinge beschreiben.
Mit den PDG-Massen $m_e = 0{,}51099895$~MeV, $m_\mu = 105{,}6583755$~MeV und
$m_\tau = 1776{,}93 \pm 0{,}09$~MeV ist $Q - 2/3 = -2{,}2\times10^{-6}$, relativ
$-0{,}00033\,\%$; die $m_\tau$-Unsicherheit ergibt $\pm 0{,}0008\,\%$. Die für externe
Darstellungen gerundete Angabe $0{,}001\,\%$ entspricht damit der Messgenauigkeit.
Die mit den T0-Massen $m_i = r_i\xi^{p_i}v$ berechnete Abweichung beträgt dagegen $+0{,}161\,\%$.
